
\documentclass{article}
\usepackage{colortbl}
\usepackage{makecell}
\usepackage{multirow}
\usepackage{supertabular}

\begin{document}

\newcounter{utterance}

\centering \large Interaction Transcript for game `eqbench', experiment `eq\_bench', episode 21 with reflexion{-}Qwen3.5{-}27B.
\vspace{24pt}

{ \footnotesize  \setcounter{utterance}{1}
\setlength{\tabcolsep}{0pt}
\begin{supertabular}{c@{$\;$}|p{.15\linewidth}@{}p{.15\linewidth}p{.15\linewidth}p{.15\linewidth}p{.15\linewidth}p{.15\linewidth}}
    \# & \multicolumn{2}{c}{Player} && \multicolumn{2}{c}{Game Master} \\
    \hline

    \theutterance \stepcounter{utterance}  
    & & & \multicolumn{4}{p{0.6\linewidth}}{
        \cellcolor[rgb]{0.9,0.9,0.9}{
            \makecell[{{p{\linewidth}}}]{
                \texttt{\tiny{[P1$\langle$GM]}}
                \\ 
\texttt{Your task is to predict the likely emotional responses of a character in this dialogue:} \\
\\ 
\texttt{Elliot: Alex, there's something I need to tell you.} \\
\texttt{Alex: Sure, man. What's up?} \\
\texttt{Elliot: I... I've been struggling with something. Something I can't keep to myself anymore.} \\
\texttt{Alex: Okay, you're starting to worry me. What is it?} \\
\texttt{Elliot: It's about... me. And my feelings. My feelings for... someone.} \\
\texttt{Alex: Well, who is it? Sophie? She's pretty cute.} \\
\texttt{Elliot: No, it's not Sophie. It's... It's you, Alex.} \\
\texttt{Alex: Me? What do you mean?} \\
\texttt{Elliot: I mean I have feelings for you. More than just friendship.} \\
\texttt{Alex: But... I'm with Sophie.} \\
\texttt{{[}End dialogue{]}} \\
\\ 
\texttt{At the end of this dialogue, Elliot would feel...} \\
\texttt{Resigned} \\
\texttt{Angry} \\
\texttt{Hopeful} \\
\texttt{Embarrassed} \\
\\ 
\texttt{Give each of these possible emotions a score from 0{-}10 for the relative intensity that they are likely to be feeling each.} \\
\\ 
\texttt{You must output in the following format, including headings (of course, you should give your own scores), with no additional commentary:} \\
\\ 
\texttt{Resigned: <score>} \\
\texttt{Angry: <score>} \\
\texttt{Hopeful: <score>} \\
\texttt{Embarrassed: <score>} \\
\\ 
\\ 
\texttt{{[}End of answer{]}} \\
\\ 
\texttt{Remember: zero is a valid score, meaning they are likely not feeling that emotion. You must score at least one emotion > 0.} \\
\\ 
\texttt{Your answer:} \\
            }
        }
    }
     \\ \\

    \theutterance \stepcounter{utterance}  
    & \multicolumn{4}{p{0.6\linewidth}}{
        \cellcolor[rgb]{0.561,0.737,0.561}{
            \makecell[{{p{\linewidth}}}]{
                \texttt{\tiny{[P1$\rangle$GM]}}
                \texttt{Resigned: 2} \\
\texttt{Angry: 1} \\
\texttt{Hopeful: 1} \\
\texttt{Embarrassed: 8} \\
            }
        }
    }
    & & \\ \\

    \theutterance \stepcounter{utterance}  
    & & & \multicolumn{2}{p{0.3\linewidth}}{
        \cellcolor[rgb]{0.9,0.9,0.9}{
            \makecell[{{p{\linewidth}}}]{
                \texttt{\tiny{[GM$|$GM]}}
                \texttt{Resigned: 2} \\
\texttt{Angry: 1} \\
\texttt{Hopeful: 1} \\
\texttt{Embarrassed: 8} \\
            }
        }
    }
    & & \\ \\

    \theutterance \stepcounter{utterance}  
    & & & \multicolumn{2}{p{0.3\linewidth}}{
        \cellcolor[rgb]{0.9,0.9,0.9}{
            \makecell[{{p{\linewidth}}}]{
                \texttt{\tiny{[GM$|$GM]}}
                \texttt{Resigned: 7} \\
\texttt{Angry: 0} \\
\texttt{Hopeful: 3} \\
\texttt{Embarrassed: 8} \\
            }
        }
    }
    & & \\ \\

    \theutterance \stepcounter{utterance}  
    & & & \multicolumn{2}{p{0.3\linewidth}}{
        \cellcolor[rgb]{0.9,0.9,0.9}{
            \makecell[{{p{\linewidth}}}]{
                \texttt{\tiny{[GM$|$GM]}}
                \texttt{game\_result = LOSE} \\
            }
        }
    }
    & & \\ \\

    \theutterance \stepcounter{utterance}  
    & & \multicolumn{4}{p{0.6\linewidth}}{
        \cellcolor[rgb]{0.95,0.95,0.95}{
            \makecell[{{p{\linewidth}}}]{
                \texttt{\tiny{[P1$\leftrightarrow$P]}}
                \texttt{Post{-}episode Reflexion (not sent to GM):} \\
\texttt{New plan: My previous attempt failed because I assigned a high score to "Embarrassed" while neglecting the more dominant emotion of "Resigned" in the face of rejection. When Elliot confesses and is immediately told Alex is with someone else, the primary reaction is the realization that his hope is futile, leading to resignation, rather than just shame. Next time, I will analyze the dialogue's conclusion to identify the most logical emotional consequence of the rejection. I will assign the highest score (e.g., 9 or 10) to "Resigned" to reflect the acceptance of the situation, give a moderate score to "Embarrassed" (e.g., 4) for the awkwardness of the confession, and keep "Angry" and "Hopeful" very low (0{-}1) as they are less likely given the immediate closure of the conversation. I will ensure the scores accurately reflect the narrative arc where the confession is met with a hard boundary.} \\
            }
        }
    }
    & \\ \\

\end{supertabular}
}

\end{document}
